\section{Ninja and his Friends (3D DP / Cherry Pickup)}
\label{sec:dp-ninja-friends}

\problemstatement{Given an $m \times n$ grid of chocolates, two
collectors start at the top row -- one at $(0,0)$, the other at
$(0,n-1)$ -- and both move \textbf{simultaneously}, one row at a time,
down to row $m-1$. From $(\textit{row},\textit{col})$, a collector may
move to $(\textit{row}+1,\textit{col}-1)$,
$(\textit{row}+1,\textit{col})$, or $(\textit{row}+1,\textit{col}+1)$.
Each collects the chocolates in every cell it visits; if both collectors
occupy the \textbf{same} cell at the same time, that cell's chocolates
are counted only \textbf{once}. Maximize the total chocolates collected.}

\begin{patternbox}
This closing problem of the chapter is a deliberate capstone: \emph{``two
agents moving simultaneously through the same grid''} is
Section~\ref{sec:dp-grids-intro}'s 3D DP keyword pattern. The state must
track \textbf{three} coordinates -- the shared row (since both move in
lockstep, one row per step) and \emph{each} collector's column -- because
the optimal move for one collector can depend on where the other
currently is (specifically, whether they currently coincide, which
determines whether a cell's chocolates are double-counted).
\end{patternbox}

\subsection*{Deriving the recurrence}

Let $f(\textit{row}, \textit{col1}, \textit{col2})$ be the maximum
chocolates collectible from row $\textit{row}$ down to the last row,
given that collector $1$ is currently at column $\textit{col1}$ and
collector $2$ is currently at column $\textit{col2}$ (both in the same
row, since they move together). The chocolates gained at this row are
$\textit{grid}[\textit{row}][\textit{col1}]$ plus
$\textit{grid}[\textit{row}][\textit{col2}]$ \emph{unless}
$\textit{col1}=\textit{col2}$, in which case the cell's value is added
only once. Each collector independently chooses one of $3$ column
moves ($-1, 0, +1$) for the next row, giving $3 \times 3 = 9$
combinations to maximize over:
\[
  f(\textit{row}, \textit{col1}, \textit{col2}) = \textit{collected} +
  \max_{\substack{d_1, d_2 \in \{-1,0,1\} \\ \text{both columns valid}}}
  f(\textit{row}+1, \textit{col1}+d_1, \textit{col2}+d_2),
\]
where $\textit{collected} = \textit{grid}[\textit{row}][\textit{col1}] +
(\textit{col1} \ne \textit{col2} ? \textit{grid}[\textit{row}][\textit{col2}]
: 0)$. Base case: at the last row ($\textit{row}=m-1$), there is no
further move, so $f(m-1,\textit{col1},\textit{col2}) =
\textit{collected}$ directly. The overall answer is $f(0, 0, n-1)$.

\subsection*{Approach 1: Recursion}

\begin{lstlisting}
#include <bits/stdc++.h>
using namespace std;

int solve(int row, int col1, int col2, vector<vector<int>>& grid, int m, int n) {
    if (col1 < 0 || col1 >= n || col2 < 0 || col2 >= n) return INT_MIN;

    int collected = grid[row][col1];
    if (col1 != col2) collected += grid[row][col2];

    if (row == m - 1) return collected;

    int best = INT_MIN;
    for (int d1 = -1; d1 <= 1; d1++) {
        for (int d2 = -1; d2 <= 1; d2++) {
            int next = solve(row + 1, col1 + d1, col2 + d2, grid, m, n);
            if (next != INT_MIN) best = max(best, next);
        }
    }
    return collected + best;
}

int main() {
    vector<vector<int>> grid = {{3,1,1},{2,5,1},{1,5,5},{2,1,1}};
    int m = grid.size(), n = grid[0].size();
    cout << solve(0, 0, n - 1, grid, m, n) << endl; // 24
    return 0;
}
\end{lstlisting}

\begin{complexitybox}[Approach 1 Complexity]
\textbf{Time.} $O(9^m)$ -- up to $9$ branches per call, depth $m$.
\textbf{Space.} $O(m)$ recursion stack.
\end{complexitybox}

\subsection*{Approach 2: Memoization}

\begin{lstlisting}
#include <bits/stdc++.h>
using namespace std;

int solveMemo(int row, int col1, int col2, vector<vector<int>>& grid,
              int m, int n, vector<vector<vector<int>>>& memo) {
    if (col1 < 0 || col1 >= n || col2 < 0 || col2 >= n) return INT_MIN;

    int collected = grid[row][col1];
    if (col1 != col2) collected += grid[row][col2];

    if (row == m - 1) return collected;
    if (memo[row][col1][col2] != -1) return memo[row][col1][col2];

    int best = INT_MIN;
    for (int d1 = -1; d1 <= 1; d1++) {
        for (int d2 = -1; d2 <= 1; d2++) {
            int next = solveMemo(row + 1, col1 + d1, col2 + d2, grid, m, n, memo);
            if (next != INT_MIN) best = max(best, next);
        }
    }
    return memo[row][col1][col2] = collected + best;
}

int main() {
    vector<vector<int>> grid = {{3,1,1},{2,5,1},{1,5,5},{2,1,1}};
    int m = grid.size(), n = grid[0].size();
    vector<vector<vector<int>>> memo(m, vector<vector<int>>(n, vector<int>(n, -1)));
    cout << solveMemo(0, 0, n - 1, grid, m, n, memo) << endl; // 24
    return 0;
}
\end{lstlisting}

\begin{complexitybox}[Approach 2 Complexity]
\textbf{Time.} $O(m \cdot n^2 \cdot 9)= O(m \cdot n^2)$ -- $m \cdot n^2$
distinct states, each doing $O(9)=O(1)$ work.
\textbf{Space.} $O(m \cdot n^2)$ memo $+$ $O(m)$ recursion stack.
\end{complexitybox}

\subsection*{Approach 3: Tabulation}

\begin{lstlisting}
#include <bits/stdc++.h>
using namespace std;

int solveTabulation(vector<vector<int>>& grid) {
    int m = grid.size(), n = grid[0].size();
    vector<vector<vector<int>>> dp(m, vector<vector<int>>(n, vector<int>(n, 0)));

    for (int col1 = 0; col1 < n; col1++) {
        for (int col2 = 0; col2 < n; col2++) {
            dp[m-1][col1][col2] = grid[m-1][col1] +
                (col1 != col2 ? grid[m-1][col2] : 0);
        }
    }

    for (int row = m - 2; row >= 0; row--) {
        for (int col1 = 0; col1 < n; col1++) {
            for (int col2 = 0; col2 < n; col2++) {
                int collected = grid[row][col1] +
                    (col1 != col2 ? grid[row][col2] : 0);

                int best = INT_MIN;
                for (int d1 = -1; d1 <= 1; d1++) {
                    for (int d2 = -1; d2 <= 1; d2++) {
                        int nc1 = col1 + d1, nc2 = col2 + d2;
                        if (nc1 >= 0 && nc1 < n && nc2 >= 0 && nc2 < n) {
                            best = max(best, dp[row+1][nc1][nc2]);
                        }
                    }
                }
                dp[row][col1][col2] = collected + (best == INT_MIN ? 0 : best);
            }
        }
    }
    return dp[0][0][n-1];
}

int main() {
    vector<vector<int>> grid = {{3,1,1},{2,5,1},{1,5,5},{2,1,1}};
    cout << solveTabulation(grid) << endl; // 24
    return 0;
}
\end{lstlisting}

\subsection*{Dry run}

\dryrun{Take $\textit{grid} = \begin{bmatrix}3&1&1\\2&5&1\\1&5&5\\2&1&1\end{bmatrix}$.}

The last row's table is filled directly: e.g.\
$\textit{dp}[3][0][2] = \textit{grid}[3][0]+\textit{grid}[3][2] =
2+1=3$ (collectors at different columns, both counted). Working
upward, $\textit{dp}[2][\cdot][\cdot]$ combines row $2$'s chocolates
with the best reachable entry of $\textit{dp}[3]$, and so on up to
$\textit{dp}[0][0][2]$. Tracing the full optimal path (confirmed by the
code): collector $1$ follows $0\to1\to2\to1$ (columns), collector $2$
follows $2\to1\to2\to1$, jointly visiting cells that sum to a maximum
of $\boxed{24}$ chocolates.

\begin{complexitybox}[Approach 3 Complexity]
\textbf{Time.} $O(m \cdot n^2)$.
\textbf{Space.} $O(m \cdot n^2)$ table, no recursion stack.
\end{complexitybox}

\subsection*{Approach 4: Space Optimization}

\begin{ideabox}[Only the Next Row's 2D Slice Is Ever Needed]
Exactly as in Section~\ref{sec:dp-unique-paths}, row $\textit{row}$'s
table only ever reads from row $\textit{row}+1$ -- so the full 3D table
collapses to two rolling 2D slices, each of size $n \times n$, instead of
$m \times n \times n$.
\end{ideabox}

\begin{lstlisting}
#include <bits/stdc++.h>
using namespace std;

int solveSpaceOptimized(vector<vector<int>>& grid) {
    int m = grid.size(), n = grid[0].size();
    vector<vector<int>> nextRow(n, vector<int>(n, 0));

    for (int col1 = 0; col1 < n; col1++)
        for (int col2 = 0; col2 < n; col2++)
            nextRow[col1][col2] = grid[m-1][col1] +
                (col1 != col2 ? grid[m-1][col2] : 0);

    for (int row = m - 2; row >= 0; row--) {
        vector<vector<int>> currentRow(n, vector<int>(n, 0));
        for (int col1 = 0; col1 < n; col1++) {
            for (int col2 = 0; col2 < n; col2++) {
                int collected = grid[row][col1] +
                    (col1 != col2 ? grid[row][col2] : 0);

                int best = INT_MIN;
                for (int d1 = -1; d1 <= 1; d1++) {
                    for (int d2 = -1; d2 <= 1; d2++) {
                        int nc1 = col1 + d1, nc2 = col2 + d2;
                        if (nc1 >= 0 && nc1 < n && nc2 >= 0 && nc2 < n) {
                            best = max(best, nextRow[nc1][nc2]);
                        }
                    }
                }
                currentRow[col1][col2] = collected + (best == INT_MIN ? 0 : best);
            }
        }
        nextRow = currentRow;
    }
    return nextRow[0][n-1];
}

int main() {
    vector<vector<int>> grid = {{3,1,1},{2,5,1},{1,5,5},{2,1,1}};
    cout << solveSpaceOptimized(grid) << endl; // 24
    return 0;
}
\end{lstlisting}

\begin{complexitybox}[Approach 4 Complexity]
\textbf{Time.} $O(m \cdot n^2)$.
\textbf{Space.} $O(n^2)$ for two rolling 2D slices, instead of
$O(m \cdot n^2)$ for the full 3D table.
\end{complexitybox}

\begin{takeawaybox}
A 3D DP state is exactly what ``$k$ agents moving through the same
structure simultaneously'' demands: one shared dimension for the
synchronized step (here, the row both collectors are always on
together), plus one dimension per agent for its own position. Everything
about the four-stage methodology -- recursion first, then memoization,
then tabulation in the base case's direction, then collapsing to rolling
slices -- is unchanged from the simpler 1D and 2D problems earlier in
this book; only the number of dimensions being tracked grows.
\end{takeawaybox}
